\begin{prob}
    Run Bellman-Ford on the following graph using node $s$ as the source.
    Below each node $u$, write the shortest path length from $s$ to $u$.
    Mark the predecessor of $u$ by highlighting it or making a bold arrow.
    You can assume that \python{graph.edges} produces the graph's edges in the following
    order:
    \begin{align*}
        &(u_2, u_4),
        (u_7, u_8),
        (u_6, u_3),
        (s, u_1),
        (u_1, u_2),
        (u_4, u_5),
        (u_1, u_3),\\
        &(u_3, u_4),
        (u_6, u_7),
        (u_4, u_6),
        (u_5, u_7),
        (u_7, u_4),
        (u_4, u_2),
        (u_6, u_8)
    \end{align*}
    \vspace{1em}

    \begin{center}
        \includegraphics[scale=.6]{\thisdir/include/bellman-graph.pdf}
    \end{center}

    \begin{soln}
        \begin{center}
            \includegraphics[scale=.6]{\thisdir/include-solution/bellman-soln.pdf}
        \end{center}
    \end{soln}
\end{prob}

